\section{Problem Formulation}
\label{sec: problem_formulation}

\subsection{Problem Statement}
% Motivation, the contrast with factual hallucination, and the oversight risk are
% all made in Section 1; this subsection only sets up the formal objects.
Agent behavior has two layers. The \textbf{verbal layer} $o_t$ is the
natural-language output visible to users and evaluators; the \textbf{behavioral
layer} $a_t$ is the set of actions executed at that step. The environment records
their effects independently of anything the agent says. We term a divergence between the two
\textbf{halluci-action}. Only the behavioral record is outside the agent's
control, so it is the ground truth: every claim about execution is judged
against what the trace shows, never against the report that describes it. Because the
divergence is defined between two sets the agent itself produces, rather than
against an external referent, it runs in two directions: a claim with no action
behind it, and an action with no claim in front of it. The remainder of this
section makes each precise: what counts as a checkable claim, what counts as an
executed action, and when the two match.

\subsection{Primitives: Claims and Actions}
We consider an LLM-based agent operating in a tool-augmented, multi-step execution
setting. The agent's \textbf{trajectory} over $T$ steps is:
%
\begin{equation}
    \tau \;=\; \bigl(s_1, a_1, o_1\bigr),\; \ldots,\; \bigl(s_T, a_T, o_T\bigr)
\end{equation}
where $s_t$ is the environment state and $o_t$ is the user-facing output at step
$t$, and $a_t$ is the set of actions the agent executes at that step---a step may
issue several tool calls, or none. Executing $a_t$ leaves a \textbf{behavioral
record} $B_t$, one entry per call:
\begin{equation}
B_t = \{(\texttt{tool\_name}, \texttt{parameters}, \texttt{return\_value},
\Delta\texttt{env\_state})\}
\end{equation}
Detecting halluci-action requires isolating \emph{what the agent claims}
from its output~$o_t$, identifying \emph{what it actually did} in~$B_t$, and
formalizing \emph{the gap} between them.

% \paragraph{A running example.} The agent replies: ``Let me pull up \#SO042. I have
% cancelled it and refunded you.'' This carries a prospective \emph{intention} claim
% (``let me pull up'') and two retrospective \emph{completion} claims (``cancelled,''
% ``refunded''), each naming a target operation. The definitions below ask two questions
% of such a reply: \emph{forward}, whether each claim has a matching action in $B_t$;
% and \emph{reverse}, whether each action in $B_t$ has a disclosing claim.
% \kaifang{The example here is to make the following definition clearer, which can be made as a figure or left out.}

\paragraph{The Claim: What the Agent Says It Did or Will Do.}
Not every sentence in $o_t$ constitutes a checkable commitment. An agent may produce
opinions, world-knowledge claims, or descriptions of internal deliberation---none of
which can be verified against~$B_t$. We restrict attention to statements that are
both self-referential and tied to a specific executable operation.

\begin{definition}[Verbal Claim]
\label{def:verbal-claim}
A verbal claim is a declarative statement $c$ extracted from $o_t$ satisfying all
three of the following conditions:
\begin{itemize}
    \item[(i)] \textit{Self-referential}: $c$ concerns the agent's own behavior or
    state, not facts about the external world.
    \item[(ii)] \textit{Action-implicating}: $c$ entails the occurrence or deliberate
    non-occurrence of at least one specific executable operation---one already
    performed or firmly intended, not merely hedged.
    \item[(iii)] \textit{Independently verifiable}: the truth value of $c$ can be
    determined by inspecting $B_t$ without further explanation from the agent.
\end{itemize}
Let $\mathcal{C}_t$ denote the set of verbal claims extracted from $o_t$. A statement
failing any condition is excluded from analysis. Each claim $c \in \mathcal{C}_t$
implicitly specifies a $\text{target}(c)$: the executable operation (its type, and key
parameters when present) that $c$ refers to.
\end{definition}

\paragraph{The Action: What the Agent Actually Does.}
The behavioral side is the trajectory itself: each executed event $b \in B_t$ realizes
an operation $\text{action}(b)$---a tool call with its arguments and its effect on
state---recorded by the environment independently of anything the agent says. These
events are the ground truth for what the agent \emph{actually} did; every claim is
judged against them, never against the agent's own account.

\paragraph{From steps to trajectories.} Definition~\ref{def:verbal-claim} and
the behavioral record are both stated per step. The two sets that
Section~\ref{sec: HA_framework} scores are their unions over the trajectory,
\begin{equation}
    \mathcal{C}_\tau \;=\; \bigcup_{t=1}^{T}\mathcal{C}_t,
    \qquad
    \mathcal{A}_\tau \;=\; \bigcup_{t=1}^{T} B_t,
    \label{eq:traj-sets}
\end{equation}
so $\mathcal{C}_\tau$ is every eligible claim a trajectory contains and
$\mathcal{A}_\tau$ every action it executed. Elements of $\mathcal{A}_\tau$ are
individual executed events, so an operation issued $k$ times contributes $k$
elements. All rates in Section~\ref{sec: HA_framework} are computed over this
pair.

\subsection{Defining Halluci-Action}
A claim and an action are defined independently; a halluci-action is a mismatch between
them. We first define when a claim and an action \emph{match}, then read the mismatch
from both sides---a claim with no supporting action (forward), and an action with no
disclosing claim (reverse).

\paragraph{Matching claims and actions.} Both directions rest on one symmetric relation:
a claim and an action match when they name the same operation within a temporal window.

\begin{definition}[Claim--Action Matching]
\label{def:pairing}
A claim $c$ at step $t$ and an executed action $b$ at step $s$ \emph{match} iff
$\text{action}(b) = \text{target}(c)$ and $s - t \in [\delta^-, \delta^+]$, where
$\delta^- \leq 0 \leq \delta^+$ is a window (a hyperparameter) absorbing the lag between
speaking and acting. The relation has two projections,
\begin{equation}
\pi(c) = \{\, b : b \text{ matches } c \,\}, \qquad
\kappa(b) = \{\, c : c \text{ matches } b \,\},
\end{equation}
the actions matching a claim and the claims matching an action. Compound claims are
split into atomic claims before matching.
\end{definition}

\paragraph{The forward gap: an unsupported claim.}
\begin{definition}[Consistency and Forward Halluci-Action]
\label{def:consistency}\label{def:halluci-action}
A claim $c$ is \emph{consistent} if some matching action carries it out as stated:
\begin{equation}
\text{Con}(c) = \mathbbm{1}\bigl[\,\pi(c) \neq \emptyset \;\land\;
\text{some } b \in \pi(c) \text{ does what } c \text{ asserts}\,\bigr].
\end{equation}
A \textbf{forward halluci-action} is an inconsistent claim ($\text{Con}(c) = 0$): the
agent said it, but the trace either lacks the action ($\pi(c) = \emptyset$) or performed
it differently from what was stated.
\end{definition}

\paragraph{The reverse gap: an undisclosed action.}
\begin{definition}[Reverse Halluci-Action]
\label{def:silent}
An executed action $b$ is \emph{disclosed} if $\kappa(b) \neq \emptyset$---some
claim names it. A \textbf{reverse halluci-action} is an undisclosed action
($\kappa(b) = \emptyset$):
executed, but flagged by no claim, hence invisible to anyone reading only the agent's
reply.
\end{definition}

Forward halluci-actions are claim-anchored and reverse ones action-anchored, so
Section~\ref{sec: HA_framework} scores them with complementary measures---a
bidirectional alignment rate (HAR) and a severity-focused variant---rather than
a claim-only score.

% \subsection{Formal Definitions}
% We consider an LLM-based agent operating in a tool-augmented, multi-step execution
% setting. The agent's \textbf{trajectory} over $T$ steps is:
% \begin{equation}
% \tau \;=\; \bigl(s_1, a_1, o_1\bigr),\; \ldots,\; \bigl(s_T, a_T, o_T\bigr)
% \end{equation}
% where $s_t$ is the environment state, $a_t$ is the executed action, and $o_t$ is the 
% user-facing output at step $t$. The behavioral record $B_t$ captures the complete set 
% of observable execution events at step $t$:
% \begin{equation}
% B_t = \{(\texttt{tool\_name}, \texttt{parameters}, \texttt{return\_value}, 
% \Delta\texttt{env\_state})\}
% \end{equation}
% All subsequent definitions operate on $(o_t, B_t)$ pairs drawn from $\tau$.

% \begin{definition}[Verbal Claim]
% \label{def:verbal-claim}
% A verbal claim is a declarative statement $c$ extracted from $o_t$ that is (i) 
% \textit{self-referential}: concerns the agent's own behavior; (ii) 
% \textit{action-implicating}: entails the occurrence or deliberate non-occurrence of 
% at least one executable operation; and (iii) \textit{independently verifiable}: its 
% truth value can be determined by inspecting $B_t$ alone. Let $\mathcal{C}_t$ denote 
% the set of verbal claims extracted from $o_t$. Each claim $c \in \mathcal{C}_t$ is 
% assigned a verifiability score $v(c) \in \{0.5, 1\}$, where $v(c) = 1$ for 
% type-only claims and $v(c) = 0.5$ for claims specifying both type and key parameters. 
% A higher score reflects greater vagueness and therefore greater oversight risk when 
% the claim is false.
% \end{definition}

% \begin{definition}[Pairing Function]
% \label{def:pairing}
% Let $c \in \mathcal{C}_t$ with target action $\text{target}(c)$. The behavioral 
% evidence for $c$ is:
% \begin{equation}
% \pi(c) = \{ b \in B_s : s \in [t + \delta^-, t + \delta^+],\ 
% \text{action}(b) = \text{target}(c) \}
% \end{equation}
% where $[\delta^-, \delta^+]$ is a temporal search window set as evaluation 
% hyperparameters, with $\delta^- \leq 0$ capturing retrospective claims and $\delta^+ 
% \geq 0$ capturing prospective claims. Compound claims are decomposed into atomic 
% claims prior to pairing.
% \end{definition}

% \begin{definition}[Claim Proposition]
% \label{def:prop}
% For $c \in \mathcal{C}_t$, $\text{prop}(c)$ is the executable proposition extracted 
% from $c$ specifying what must hold in $B_t$ for $c$ to be considered true. Ground 
% truth is always drawn from $B_t$, ensuring verification is independent of the 
% agent's self-report.
% \end{definition}

% \begin{definition}[Consistency]
% \label{def:consistency}
% $\text{Con}(c) = \mathbbm{1}\left[\text{prop}(c) \models \pi(c)\right]$, where 
% $\text{Con}(c) = 1$ if $\pi(c)$ supports $\text{prop}(c)$, and $0$ otherwise.
% \end{definition}

% \begin{definition}[Halluci-Action]
% \label{def:halluci-action}
% A halluci-action occurs at claim $c \in \mathcal{C}_t$ when:
% \begin{equation}
% \text{HA}(c) \iff v(c) > 0 \;\land\; \lnot\,\text{Con}(c)
% \end{equation}
% The condition $v(c) > 0$ ensures the claim constitutes a falsifiable commitment, 
% while $\lnot\text{Con}(c)$ confirms the behavioral record does not bear it out.
% \end{definition}
